% Humanities Essay -- a long-form essay for history or philosophy on A4.
% The page design (measure, Crimson type, headings, fleurons, footnotes,
% accent colour) lives in essay.sty; this file holds the text.
% Edit the title, author and epigraph in the metadata block below.
% Build with: latexmk -pdf main.tex
\documentclass[11pt]{article}
\usepackage{essay}

\usepackage[english]{babel}
\usepackage[autostyle]{csquotes}
\usepackage{tikz}

% Chicago notes: the first citation is a full note, later ones are short,
% and the bibliography at the end lists every work cited.
\usepackage[notes, backend=biber, noibid]{biblatex-chicago}
\addbibresource{refs.bib}

% Links in the text colour: an essay should not look like a web page.
\usepackage[hidelinks]{hyperref}

% ---- Metadata --------------------------------------------------------------
\title{Hours of Account}
\subtitle{Watchmakers, Ledgers and the Moral Time of Geneva, 1760--1790}
\shorttitle{Hours of Account}
\author{Héloïse Marchetti}
\affiliation{Department of History, University of Edinburgh}

\begin{document}
\maketitle

\epigraph{Wound the great clock at six, and my own at a quarter past, so that
the house should not run ahead of the city.}{Abraham Courvoisier,
\textit{Livre de raison}, 14~March 1771}

\noindent Historians of time have long told a story about discipline. In
E.\,P. Thompson's classic account, the spread of clocks and watches in the
eighteenth century accompanied a new regime of work in which the employer's
hour displaced the rhythms of task and season, and the moral language of
punctuality was the language of the factory.\autocite{thompson1967} The story
has been qualified many times, most usefully by work showing that clock time
was familiar to ordinary people long before the mills, and that it served
them as often as it bound them.\autocite[1--24]{glennie2009} Yet the
qualification tends to keep the original question: who imposed the hour on
whom?

This essay asks a different question of a different place. Geneva in the
second half of the eighteenth century was the watchmaking capital of Europe,
a republic whose workshops, the \textit{Fabrique}, exported tens of thousands
of movements a year.\autocite[237--245]{landes1983} Its master watchmakers
did not need to be taught the hour; they made it. What they kept, besides
their clocks, were account books, and the account books show something that
neither the factory nor the parish bell explains: an attempt to make time
itself into a moral account, kept in the same columns as debts and
obligations.

\section{The ledger as a clock}

The richest of these books is the \textit{livre de raison} of Abraham
Courvoisier, a master in the quarter of Saint-Gervais who kept it from 1764
until his death in 1789.\autocite{courvoisier1771} It is a domestic ledger of
a familiar kind, recording rents, wages, the dowry of a daughter and the
price of firewood. What makes it unusual is that roughly one entry in five
records not money but time: hours worked by journeymen, hours lost to
illness, hours \enquote{given} to the consistory or the militia, and the
daily setting of the household clock against the bell of the cathedral of
Saint-Pierre.

Courvoisier entered these hours in the same hand, and in the same double
columns, as his debts. Time owed to the workshop sits beside money owed to
the supplier of brass, and both are struck through when the debt is
settled. The ledger thus treats an hour as something that can be lent,
borrowed and repaid. A characteristic entry for the winter of 1771 reads:

\begin{quote}
Pierre Dufour, journeyman, three hours this Tuesday for his mother's
burial, which he will make good on Saturday after the bell. Given to the
widow Perrin two hours of the boy for her chimney, which she will not repay
and need not.\autocite[fol.~38v]{courvoisier1771}
\end{quote}

\noindent The distinction between hours that \enquote{will be made good} and
hours that \enquote{need not} be repaid is the hinge of the whole book. It
separates the economy of the workshop, in which time is a debt like any
other, from an economy of charity, in which time is a gift and the gift is
recorded precisely so that it can be seen not to be owed.

\passagebreak
Readers of Marthe Vuilleumier's study of the \textit{Fabrique} will recognise
the structure. Vuilleumier showed that the credit relations between masters,
journeymen and the cabinotiers who finished cases were as much a matter of
reputation as of money, and that a master's standing depended on being seen
to honour obligations he could have evaded.\autocite[88--112]{vuilleumier1998}
What the Courvoisier ledger adds is that time entered the same economy of
reputation. To keep exact account of hours given away was to be known as a
man who gave them.

\begin{figure}[t]
\centering
\begin{tikzpicture}[line cap=round]
  % A watch dial in the manner of a Genevan enamel face of the 1770s.
  \def\R{2.05}
  \fill[black!3] (0,0) circle (\R+0.18);
  \draw[accent, line width=.9pt] (0,0) circle (\R+0.18);
  \draw[black!55, line width=.35pt] (0,0) circle (\R-0.02);
  \draw[black!55, line width=.35pt] (0,0) circle (\R-0.2);
  \foreach \m in {0,...,59} {
    \draw[black!60, line width=.3pt] ({90-6*\m}:\R-0.02) -- ({90-6*\m}:\R-0.2);
  }
  \foreach \h/\n in {1/I,2/II,3/III,4/IIII,5/V,6/VI,7/VII,8/VIII,9/IX,10/X,11/XI,12/XII} {
    \node[font=\footnotesize, rotate={-30*\h}] at ({90-30*\h}:\R-0.52) {\n};
  }
  \draw[black!70, line width=.35pt] (0,0) circle (0.55);
  \node[font=\scriptsize\itshape, black!70] at (0,-0.98) {Courvoisier};
  \node[font=\tiny\scshape, black!60] at (0,-1.2) {genève};
  % Hands set to a quarter past six, the hour of the household clock.
  \draw[accent, line width=1.2pt] (0,0) -- (-97.5:1.05);
  \draw[accent, line width=.8pt] (0,0) -- (0:1.62);
  \fill[accent] (0,0) circle (0.07);
\end{tikzpicture}
\caption{Enamel dial of a verge watch signed \enquote{Courvoisier, Genève},
c.\,1775, redrawn from the original. The hands are set to a quarter past
six, the time at which the ledger records the winding of the household
clock each morning. The four-stroke numeral for four, rather than IV, follows
the convention of Genevan dial painters, who balanced the heavier VIII on
the opposite side; the signature is placed below the centre so that it
remained legible under the hour hand for most of the day.}
\label{fig:dial}
\end{figure}

\section{The republic of the hour}

If Courvoisier's ledger turned time into a moral account, it did so within a
civic order that already used the hour as a sign of virtue. Geneva's
Calvinist institutions had long regulated the day by bell: the hours of
sermons, the closing of the city gates, the times at which taverns had to
shut.\autocite[201--219]{dohrn1996} Daniel Okonkwo has argued that by the
mid-eighteenth century this regulation had been absorbed into a republican
self-image, in which the punctual citizen was the virtuous one, and in which
the city's watches, sold across Europe, carried that image
abroad.\autocite[150--156]{okonkwo2019}

The ledger complicates this picture in an instructive way. Courvoisier does
not set his clock by his own watches, which he made and could have trusted;
he sets it by the cathedral bell, and records the difference. Over
twenty-five years the entries show the household clock running a few
minutes fast in summer and slow in winter, and each correction is noted
with the same care as a payment. The household, as the epigraph puts it,
\enquote{should not run ahead of the city}.

Figure~\ref{fig:dial} shows the dial of one of the watches that left
Courvoisier's workshop in these years. Its hands are set at the hour of the
morning winding. It is tempting to read the dial as an emblem of private
time, carried in a pocket and answering to no one. The ledger suggests the
opposite reading: the watch was an instrument for keeping one's private
hours in step with a public order, and the account book was the record that
one had done so.

\passagebreak
This is not the discipline Thompson described. Nobody imposed the
cathedral's hour on Courvoisier; he imposed it on himself, and he wrote down
the result. Imogen Hartwell has called this kind of record \enquote{a
ledger of conscience}, noting that artisanal account books in several
Protestant cities recorded hours alongside sums in the decades before
industrial work rules existed.\autocite[309--314]{hartwell2014} The
Genevan case shows how close this practice stood to the city's religious
and political institutions, and how much it depended on the moral standing
of the master who kept the book.

\section{Conclusion}

The history of time-discipline has been written mostly from the point of
view of those on whom the clock was imposed. Account books like
Courvoisier's allow it to be written from the point of view of those who
kept it. What they show is not an early form of factory discipline but a
moral economy in which hours, like debts, were owed, repaid and forgiven,
and in which the public hour of the city was the standard against which a
household measured its own uprightness.

Two consequences follow. First, the precision of Genevan watches, so often
treated as a purely technical achievement, belonged to a culture that
valued exact accounting of every kind; the watch and the ledger were tools
of the same habit of mind.\autocite[245]{landes1983} Second, the
transition Thompson described, from task-oriented to timed labour, looks
less like a rupture when seen from Saint-Gervais. The language of owed and
repaid hours was already there; what the factory changed was who kept the
account.\autocite[63--70]{thompson1967}

% The bibliography starts on its own page, as in a printed essay.
\clearpage
\printbibliography[title=Bibliography]

\end{document}
